\section{Mass decay and the enclosure radius}
\label{sec:rpi:flow-estimates}

Area preservation and Proposition~\ref{prop:rpi:mass-evolution} reduce
the argument to the decay of $\mu(t)=m(t)-c_t$ and the location of the
frontier. The next estimates provide a fixed harmonic tail in normalized
coordinates, on which the quantitative positive-mass theorem will apply.
Put $d=n-2$ and $q=2(n-1)/d$, and retain
$\omega=|S^{n-1}|$ for the Euclidean area of the unit sphere.

\begin{proposition}[Vanishing deficit and bounded normalized radius]
\label{prop:rpi:flow-estimates}
For the flow in Proposition~\ref{prop:rpi:mass-evolution},
$\mu(t)\to0$ as $t\to\infty$. There is a fixed $R_{\max}$ such that
\begin{equation}\label{eq:rpi:radius-bound}
 \Sigma_t\subset B_{R_{\max}e^{2t/d}}\qquad(t\ge0).
\end{equation}
Here the radius is measured in the original end coordinates, with the
compact core included in the ball. Equivalently, every normalized
frontier lies in $B_{R_{\max}}$.
\end{proposition}

\begin{proof}
\textbf{The integrated mass identity.}\quad
Since $B'=b$ and $B(0)=0$,
\[
 \frac{d}{dt}(e^{-t}B(t))=e^{-t}(b(t)-B(t)).
\]
Consequently
\[
 \int_0^t\mu(s)\,ds
 =\frac12(1-e^{-2t})m(0)-e^{-t}B(t)
 =\frac{m(0)-m(t)}2.
\]
This supplies integrability of the nonnegative
$\mu$. Set $a=m-2\mu=2c_t-m$.
Nesting and the maximum principle make $e^tv_t$ nondecreasing:
the later potential is zero on the larger filled set and both have
limit $-1$ at infinity. Its end coefficient $e^tb(t)$ is therefore
nondecreasing. Since
\[
 e^{2t}a(t)=2e^tb(t)-m(0),
\]
the same holds for $e^{2t}a(t)$. This implies, distributionally,
$Da+2a\,dt\geq0$. Substitution of the integral identity gives
monotonicity of $3m-2\mu+2\int_0^t m$ (up to a constant).
Thus for $h>0$,
\[
 \mu(t+h)-\mu(t)
 \leq\frac32\bigl(m(t+h)-m(t)\bigr)
       +\int_t^{t+h}m(s)\,ds
 \leq h m(0).
\]
This one-sided slope bound and integrability imply
$\mu(t)\to0$: a value at least $\varepsilon>0$ forces an
interval immediately before it on which the integral is bounded below
by a positive constant depending only on $\varepsilon,m(0)$.
If $m(0)=0$, nonnegativity and mass monotonicity give $m=0$;
the same slope bound excludes positive point values of $\mu$.

\paragraph{The escaping-point ball.}
Let $a=e^{2/d}$ and $R(s)=R_{\max}e^{2s/d}$. Suppose
that containment holds for $0\le s\le T$ but there is a point
$p\in\Sigma_{T+1}$ with $|p|\geq aR(T)$. Use the ball
\[
 \rho=\frac{a-1}{2}R(T),\qquad
 \overline{B_\rho(p)}\subset\{|y|>R(T)\}.
\]
Here the balls are Euclidean balls in the original end coordinates.
Choose $R_{\max}$ so the indicated exterior avoids the initial obstacle
and $\tfrac12\delta\leq g_0\leq2\delta$ there.

To estimate the metric on this ball, use its total harmonic conformal factor.
Write $g_0=H^{4/d}\delta$ on the initial end, where $\Delta_\delta H=0$
and $H\to1$. Then $Hv_s$ is Euclidean harmonic outside $\Sigma_s$.
If $\Sigma_s\subset B_{R(s)}$, its boundary value on $S_{R(s)}$
is nonpositive and its limit at infinity is $-e^{-s}$, so
\[
 Hv_s(y)\leq e^{-s}\left[\left(\frac{R(s)}{|y|}\right)^d-1\right]
 \qquad (|y|\geq R(s)).
\]
Using $v_s\leq0$ on the remaining time interval gives, for $r=|y|\geq R(T)$,
\[
 H u_{T+1}\leq H-1+e^{-T}+(e^T-1)R_{\max}^d r^{-d}.
\]
Choose $H-1\leq C_0r^{-d}$ and $R_{\max}^d\geq C_0$.
The last expression is at most $2e^{-T}$. Taking also $H\geq1/2$
gives the concrete bounds
\[
 e^{-T-1}\leq u_{T+1}\leq K e^{-T},\qquad K=4,
 \qquad r\geq R(T).
\]

Every finite-perimeter variation supported in $B_\rho(p)$ remains an
admissible enclosure. Comparison of area elements therefore makes
$\Sigma_{T+1}$ Euclidean perimeter-quasiminimizing there, with ratio
\[
 \gamma=2^{n-1}(eK)^q.
\]
The isoperimetric proof of the density estimate, using the competitors
that fill or delete a ball, applies at every measure-theoretic boundary
point, including a singular point. It yields
\[
 \mathcal H_\delta^{n-1}(\Sigma_{T+1}\cap B_\rho(p))
 \geq c(n,\gamma)\rho^{n-1}.
\]
It follows that the fixed area satisfies
\[
 \begin{split}
 A&\geq 2^{-(n-1)/2}e^{-q(T+1)}
       c(n,\gamma)\rho^{n-1}\\
  &=2^{-(n-1)/2}c(n,\gamma)e^{-q}
       \left(\frac{e^{2/d}-1}{2}\right)^{n-1}R_{\max}^{n-1}.
 \end{split}
\]
All time factors cancel. The constants are independent of $T$ and of
further enlargement of $R_{\max}$, giving the required contradiction.
To cover all times, enlarge $R_{\max}$ to obtain containment on $[0,1]$
from the finite-time support bound. If containment holds on
$[0,N]$, apply the argument at each $t\in[N,N+1]$ with $T=t-1$.
This proves Equation~\eqref{eq:rpi:radius-bound} at every real time.
\end{proof}
